\documentclass[10pt,a4paper,twocolumn]{article}
\usepackage[utf8]{inputenc}
\usepackage[T1]{fontenc}
\usepackage[spanish,es-tabla]{babel}
\usepackage[margin=1.8cm]{geometry}
\usepackage{lmodern,microtype,booktabs,tabularx,array,amsmath,graphicx,xcolor,tikz,pgfplots}
\usepackage{amssymb,alltt}
\usetikzlibrary{arrows.meta,positioning}
\pgfplotsset{compat=1.18}
\usepackage[font=small,labelfont=bf]{caption}
\usepackage{titlesec,fancyhdr}
\usepackage[colorlinks=true,urlcolor=blue!50!black,citecolor=blue!50!black,linkcolor=blue!50!black]{hyperref}
\setlength{\columnsep}{0.65cm}
\setlength{\parindent}{1em}
\setlength{\parskip}{2pt}
\setlength{\headheight}{14pt}
\setlength{\emergencystretch}{1.5em}
\pagestyle{fancy}\fancyhf{}\fancyhead[L]{\small EvidenCite | Informe de avance}\fancyhead[R]{\small MAIA}\fancyfoot[C]{\thepage}
\titleformat{\section}{\large\bfseries}{\thesection.}{0.45em}{}
\titleformat{\subsection}{\normalsize\bfseries}{\thesubsection.}{0.4em}{}
\titlespacing*{\section}{0pt}{9pt}{4pt}
\titlespacing*{\subsection}{0pt}{7pt}{3pt}
\newcolumntype{Y}{>{\raggedright\arraybackslash}X}
\begin{document}
\raggedbottom
\twocolumn[\begin{center}
{\LARGE\bfseries EvidenCite: recomendación local de citas y clasificación de su función en artículos académicos}\par\vspace{6pt}
Daniel Ricardo Gallo Acevedo · Dubán Leandro Jiménez Puerta\\
Julián Andrés Malaver Montoya · Javier Mauricio Torres Ortega\\
Maestría en Inteligencia Artificial, Universidad de los Andes\\
\small Informe de avance
\end{center}\vspace{7pt}]

\section{Descripción breve de la propuesta de solución}
EvidenCite estudia si la evidencia localizada en el documento citado permite clasificar mejor la función de una cita que una representación global de ese documento. Se comparan dos entradas: \textbf{A}, contexto citante más título y resumen del citado, y \textbf{B}, el mismo contexto más hasta tres fragmentos recuperados del citado. La recuperación identifica pasajes relacionados; la clasificación asigna una función discursiva. El diseño permite estudiar tanto la predicción de la función como el aporte de la información recuperada.

Este avance presenta la construcción del conjunto pareado, su evolución de \textbf{3.232 a 8.074 instancias} y el diseño comparativo de TF-IDF, SciBERT y Qwen+LoRA. La meta sigue siendo 18.000 instancias, con 2.000 por función. Los resultados corresponden a las particiones del conjunto inicial; el conjunto ampliado incorpora la taxonomía objetivo y constituye la base del siguiente ciclo experimental.

\section{Recolección y preparación de los datos}
El proceso de recolección, estructuración y conformación del dataset se diseñó como un pipeline desacoplado en cuatro capas incrementales y secuenciales (ingesta cruda, corpus estructurado, procesamiento de párrafos y ensamble pareado), tal como se ilustra en la Figura~\ref{fig:pipeline_datos}. Este diseño ofrece ventajas decisivas para la investigación: aísla el almacenamiento de fuentes masivas inmutables de la lógica de experimentación, asegura reproducibilidad e idempotencia ante reprocesamientos, y permite la maduración progresiva del corpus. Así, se habilitó primero un conjunto preliminar de 3.232 instancias pareadas para poner a prueba el protocolo comparativo A/B, se escaló posteriormente a 8.074 instancias incorporando la taxonomía canónica mediante minado y resolución bibliográfica dirigida, y actualmente se avanza hacia la meta balanceada de 18.000 pares mediante supervisión débil, todo ello sin recomputar capas precedentes y garantizando esquemas consistentes para el modelamiento.

\begin{figure}[t]
\centering
\resizebox{\columnwidth}{!}{%
\begin{tikzpicture}[
  box/.style={
    draw=black!75, rounded corners=2pt, fill=black!3,
    align=center, font=\footnotesize, inner sep=4pt
  },
  wide/.style={box, text width=8.6cm, minimum height=9mm},
  half/.style={box, text width=4.1cm, minimum height=9mm},
  arrow/.style={-{Latex[length=2.2mm]}, semithick},
  line/.style={semithick}
]

% Nivel 1
\node[wide, anchor=north] (raw) at (0, 0) {%
  \textbf{1. Ingesta cruda}\\
  unarXive (\LaTeX{}), MultiCite y fuentes bibliográficas%
};

% Nivel 2
\node[wide, anchor=north] (struct) at (0, -1.45) {%
  \textbf{2. Corpus estructurado}\\
  Textos completos, jerarquía de secciones y oraciones%
};

% Nivel 3
\node[half, anchor=north] (p1) at (-2.25, -3.15) {%
  \textbf{3a. Contextos citantes}\\[2pt]
  Ventana ($m=3$) con cita%
};
\node[half, anchor=north] (p2) at ( 2.25, -3.15) {%
  \textbf{3b. Fragmentos citados}\\[2pt]
  Sección y cuerpo ($\le 110$ pal.)%
};

% Nivel 4
\node[wide, anchor=north] (data) at (0, -4.75) {%
  \textbf{4. Ensamble pareado y partición}\\
  Paridad estricta 1:1 y división agrupada por citante%
};

% Nivel 5
\node[half, anchor=north] (cfgA) at (-2.25, -6.45) {%
  \textbf{Configuración A}\\[2pt]
  Contexto + Título y Resumen%
};
\node[half, anchor=north] (cfgB) at ( 2.25, -6.45) {%
  \textbf{Configuración B}\\[2pt]
  Contexto + Top-3 Chunks%
};

% --- Conexiones ---

% 1 -> 2
\draw[arrow] (raw.south) -- (struct.north);

% 2 -> 3a, 3b
\path (struct.south) -- ++(0,-0.45) coordinate (mid23);
\draw[line]  (struct.south) -- (mid23);
\draw[arrow] (mid23) -| (p1.north);
\draw[arrow] (mid23) -| (p2.north);

% 3a, 3b -> 4
\draw[arrow] (p1.south) -- (data.north -| p1.south);
\draw[arrow] (p2.south) -- (data.north -| p2.south);

% 4 -> Config A, Config B
\path (data.south) -- ++(0,-0.45) coordinate (mid45);
\draw[line]  (data.south) -- (mid45);
\draw[arrow] (mid45) -| (cfgA.north);
\draw[arrow] (mid45) -| (cfgB.north);

\end{tikzpicture}%
}
\caption{Arquitectura en cuatro capas desacopladas del pipeline de datos. Garantiza trazabilidad de procedencia, paridad estricta 1:1 y generación desacoplada de las entradas para las Configuraciones A y B.}\label{fig:pipeline_datos}
\end{figure}

\subsection{De las fuentes al conjunto comparable}
El corpus se construyó combinando dos fuentes complementarias. MultiCite aporta contextos de citación en el área de procesamiento de lenguaje natural con funciones discursivas anotadas por revisores humanos \cite{Lauscher2022}. Por su parte, unarXive suministra artículos científicos estructurados a partir de fuentes \LaTeX{}, cuya red de citación interna y referencias bibliográficas permite identificar de forma determinista tanto el documento citante como el documento citado  \cite{Besrour2026}.

El procesamiento inicial difiere según el origen: mientras unarXive requiere deserializar árboles de secciones, oraciones y marcadores en particiones comprimidas, MultiCite exige una ingesta dirigida con resolución de metadatos bibliográficos. No obstante, ambos flujos convergen en un esquema común en la capa de documentos estructurados (Figura~\ref{fig:pipeline_datos}). Sobre unarXive se implementó un minado dirigido mediante reglas léxicas de alta precisión pero bajo recall, orientado no solo a obtener ejemplos adicionales, sino a brindar una asignación preliminar confiable para clases de la taxonomía canónica ausentes en MultiCite (\textit{Evidence}, \textit{Originator}, \textit{Further Reading} y \textit{Basis}). Cada instancia conserva metadatos exhaustivos de procedencia (fuente original, etiqueta heredada y mecanismo de asignación), garantizando la trazabilidad estricta del minado y diferenciando con claridad las anotaciones de oro humano de las obtenidas mediante reglas o procedimientos automáticos.

Para comparar de forma rigurosa las entradas A y B, cada instancia reúne el contexto de la cita y la información correspondiente del trabajo citado: título, resumen y fragmentos de su cuerpo. Se preservaron los encabezados originales de sección para mantener el encuadre pragmático del texto, fundamentado en la teoría de dispositivos de encuadre (\textit{framing devices}) de CiteFusion \cite{Paolini2025}, la cual demuestra que la estructura retórica autoral correlaciona fuertemente con la intención de citación frente a normalizaciones estáticas IMRaD. En casos de subsecciones anidadas o párrafos huérfanos, se aplicó el principio de ámbito envolvente más próximo (\textit{nearest enclosing scope}) para asociar deterministamente el texto al encabezado jerárquico inmediatamente superior en el árbol del documento. Asimismo, siguiendo a Lauscher et al. \cite{Lauscher2022}, se adoptó una ventana simétrica de tres oraciones ($m=3$) centrada en el marcador de citación, ajustada en los bordes del documento; esta ventana se encuentra parametrizada en la configuración del pipeline, facilitando variaciones experimentales futuras sin recomputar las capas base.

\subsection{Resolución bibliográfica, escala y paridad A/B}
La resolución bibliográfica vincula cada cita con el artículo citado correspondiente mediante un protocolo determinista de alta fidelidad: se normalizan los identificadores de arXiv tanto en su nomenclatura moderna como clásica, se aprovecha la red de enlaces internos entre marcadores y referencias en unarXive \cite{Besrour2026} y, para las citas de MultiCite, se consulta la API Graph de Semantic Scholar mediante búsquedas exactas por título con normalización de caracteres y cotejo del primer autor. Para optimizar el ancho de banda y garantizar reproducibilidad, las respuestas se almacenan en una base de datos local persistente SQLite con políticas de reintento exponencial con variación aleatoria (\textit{jitter}), evitando bloqueos por límite de tasa.

La escala del procesamiento motivó decisiones clave de implementación: se gestionaron más de 35~GB de datos en la capa de ingesta cruda y se estructuraron 69.331 documentos completos en la capa de corpus estructurado (Figura~\ref{fig:pipeline_datos}). La lectura y transformación se ejecutan mediante flujos por líneas (\textit{streaming}), asegurando un consumo de memoria acotado sin cargar el volumen masivo en memoria principal.

El principio metodológico central en la construcción del dataset es la \textbf{paridad estricta 1:1 de instancias}:
\[
 \mathcal{D}_{\text{eval}}(\text{Configuración A}) \equiv \mathcal{D}_{\text{eval}}(\text{Configuración B})
\]
Toda instancia evaluada en la Configuración A (título y resumen) debe contar obligatoriamente con su contraparte idéntica en la Configuración B (fragmentos discursivos recuperados densamente del cuerpo del trabajo citado). Si un artículo citado carece de texto completo estructurado en el corpus, se descarta inmediatamente del ensamble. Esta decisión es indispensable para que los resultados del modelamiento sean técnicamente comparables: admitir instancias en A por disponer únicamente de resumen, mientras B evalúa una fracción con texto completo, sesgaría la población experimental y restaría validez científica a los hallazgos.

Esta exigencia de paridad explica la complejidad del embudo de selección: de los 56.778 contextos de citación candidatos minados inicialmente, solo aquellos cuyos citados contaban con texto completo estructurado pudieron admitirse para la extracción de fragmentos. La versión experimental inicial reunió \textbf{3.232 instancias completas}; la ampliación posterior, sustentada en la descarga dirigida de artículos citados faltantes mediante un resolutor de arXiv y la expansión de la caché bibliográfica, elevó el corpus a \textbf{8.074 instancias pareadas} (Tabla~\ref{tab:evolucion}).

\begin{table}[t]
\centering\small
\caption{Evolución del conjunto experimental hacia la meta de 18.000 instancias.}\label{tab:evolucion}
\begin{tabular}{lrr}\toprule
Indicador & Inicial & Ampliado\\\midrule
Instancias pareadas & 3.232 & 8.074\\
Artículos citantes & 657 & 3.244\\
Origen MultiCite (oro humano) & 2.735 & 3.629\\
Origen minado unarXive & 497 & 4.445\\
Mapeo unívoco (\texttt{exact}) & 2.899 & 7.607\\
Mapeo compuesto (\texttt{partial}) & 333 & 467\\
Menos de tres fragmentos & 32 & 40\\\bottomrule
\end{tabular}
\end{table}

La verificación de los archivos confirma que todas las instancias de ambos conjuntos contienen contexto de cita, título, resumen y hasta tres fragmentos densamente recuperados ($k=3$), garantizando un mínimo estricto de al menos un fragmento para documentos atípicamente breves (únicamente 32 casos en el conjunto inicial y 40 en el ampliado presentan menos de tres fragmentos). Esta política conserva ejemplos utilizables y mantiene la paridad A/B. El conjunto creció 2,50 veces y alcanza el 44,9\% de la meta total; la distribución por función orienta las siguientes incorporaciones.

\subsection{Homologación y tratamiento de funciones múltiples}
La homologación hacia las nueve funciones objetivo se fundamenta en un análisis ontológico que contrasta las definiciones de anotación humana de MultiCite \cite{Lauscher2022} con las nueve motivaciones fundamentales decantadas en la meta-síntesis de Lyu et al. \cite{Lyu2021}. La Tabla~\ref{tab:mapeo} resume las correspondencias operativas. Este mapeo desacopla conceptualmente \textit{Basis} de \textit{Extends} para distinguir la adopción de un cimiento teórico-conceptual frente a la modificación técnica activa (\textit{Improvement}), y subsume \textit{Motivation} y \textit{Future Work} bajo \textit{Gap}, dado que ambas expresan delimitación de fronteras de investigación y necesidades no resueltas.

\begin{table}[t]
\centering\small
\caption{Reglas operativas de homologación hacia la taxonomía canónica \cite{Lyu2021}.}\label{tab:mapeo}
\begin{tabularx}{\columnwidth}{@{}>{\raggedright\arraybackslash}p{2.6cm}Y@{}}\toprule
Origen (MultiCite / minado) & Destino y criterio ontológico\\\midrule
Background & Background: contextualización temática y estado del arte.\\
Uses & Application: empleo directo de métodos o recursos.\\
Extends & Improvement: modificación, optimización o extensión.\\
Similarities / Differences & Comparison: contraste empírico o metodológico.\\
Motivation / Future Work & Gap: vacíos de investigación y limitaciones abiertas.\\
Basis (minada) & Basis: fundamento teórico-conceptual previo.\\
Evidence (minada) & Evidence: corroboración empírica y sustento de hallazgos.\\
Originator (minada) & Originator: atribución de autoría pionera seminal.\\
Further Reading (minada) & Further Reading: remisión a lecturas suplementarias.\\\bottomrule
\end{tabularx}
\end{table}

El tratamiento de citas con anotaciones múltiples en MultiCite adopta una regla posicional determinista: se normalizan las etiquetas preservando su orden original y se selecciona la primera para configurar la tarea de clasificación multiclase disyunta. La secuencia multietiqueta completa se conserva intacta en los metadatos y la instancia se marca con la condición \texttt{partial}. Por ejemplo, una anotación con \texttt{uses similarities} se proyecta a Application, mientras \texttt{similarities uses} se asigna a Comparison. Esta política garantiza una línea base reproducible sin imponer prioridades heurísticas arbitrarias y preserva la trazabilidad para auditorías contextuales.

Las Tablas~\ref{tab:3k} y~\ref{tab:8k} ilustran el comportamiento de esta homologación. En términos relativos, la proporción global de instancias con mapeo compuesto (\texttt{partial}) disminuye de 10,3\% en el conjunto inicial a 5,8\% en el ampliado. En términos absolutos, no obstante, estos casos aumentan de 333 a 467 debido a la incorporación de 894 nuevas citas de MultiCite rescatadas mediante la ampliación de la búsqueda bibliográfica y la descarga de textos completos (pasando de 2.735 a 3.629). Dado que la multietiqueta es intrínseca a la anotación humana en MultiCite ($\sim 12,5\%$), rescatar más citas de dicha fuente incrementa naturalmente los casos compuestos, mientras que las 4.445 instancias minadas son estrictamente unívocas (\texttt{exact}). La marca \texttt{exact} refleja un mapeo directo de esquema, no certeza de anotación.

\begin{table}[t]
\centering\small
\caption{Conjunto inicial: marcas de homologación por clase, antes del cambio a la taxonomía objetivo.}\label{tab:3k}
\begin{tabular}{lrrr}\toprule
Clase & Exact & Partial & Total\\\midrule
Background & 831 & 43 & 874\\
Motivation & 113 & 29 & 142\\
Uses & 828 & 67 & 895\\
Comparison & 522 & 150 & 672\\
Extends & 155 & 43 & 198\\
Future Work & 9 & 1 & 10\\
Further Reading & 259 & 0 & 259\\
Evidence & 156 & 0 & 156\\
Originator & 26 & 0 & 26\\\midrule
Total & 2.899 & 333 & 3.232\\\bottomrule
\end{tabular}
\end{table}

\begin{table}[t]
\centering\small
\caption{Conjunto ampliado: taxonomía objetivo y déficit frente a 2.000 instancias por clase.}\label{tab:8k}
\setlength{\tabcolsep}{3pt}
\begin{tabular}{lrrrr}\toprule
Clase & Exact & Partial & Total & Déficit\\\midrule
Background & 1.109 & 61 & 1.170 & 830\\
Gap & 168 & 44 & 212 & 1.788\\
Basis & 472 & 0 & 472 & 1.528\\
Comparison & 674 & 202 & 876 & 1.124\\
Application & 1.068 & 89 & 1.157 & 843\\
Improvement & 143 & 71 & 214 & 1.786\\
Evidence & 1.384 & 0 & 1.384 & 616\\
Originator & 344 & 0 & 344 & 1.656\\
Further Reading & 2.245 & 0 & 2.245 & 0\\\midrule
Total & 7.607 & 467 & 8.074 & 10.171\\\bottomrule
\end{tabular}
\end{table}

En el conjunto ampliado, 3.629 ejemplos proceden de MultiCite y 4.445 del minado: 44,9\% y 55,1\%, respectivamente. La integración de fuentes amplía la cobertura funcional: Basis, Evidence, Originator y Further Reading se obtienen mediante minado, mientras las otras cinco funciones proceden de MultiCite. Conservar esta distinción permite evaluar posteriormente el rendimiento por fuente y planear una revisión de etiquetas que cubra ambos procedimientos.

\subsection{Particiones y diversidad documental}
Las particiones reportadas en los experimentos iniciales (3.232 instancias) se agrupan por artículo citante para prevenir la fuga de información (\textit{data leakage}). Citas procedentes de un mismo manuscrito comparten autoría, terminología específica y giros retóricos que el modelo podría memorizar durante el entrenamiento, sobreestimando su rendimiento \cite{Lauscher2022,Jurgens2018}. La división estratificada por documento citante garantiza una separación estricta:
\[
 \text{Citantes}(\mathcal{D}_{\text{train}}) \cap \text{Citantes}(\mathcal{D}_{\text{dev}}) \cap \text{Citantes}(\mathcal{D}_{\text{test}}) = \emptyset,
\]
asegurando que la evaluación mida la generalización real ante manuscritos no observados (Tabla~\ref{tab:splits}).

\begin{table}[t]
\centering\small
\caption{Particiones empleadas por los experimentos reportados.}\label{tab:splits}
\begin{tabular}{lrrr}\toprule
Partición & Instancias & Citantes & Inst./citante\\\midrule
Train & 2.255 & 615 & 3,67\\
Dev & 491 & 26 & 18,88\\
Test & 486 & 16 & 30,38\\\bottomrule
\end{tabular}
\end{table}

Esta estricta partición documental conllevó una menor diversidad en la evaluación inicial de 3.232 instancias: cinco artículos citantes concentran el 41,2\% de la prueba y diez abarcan el 70,4\%. Asimismo, el conjunto de prueba original contiene únicamente clases provenientes de MultiCite (Background, Uses, Comparison, Extends y Motivation). La ausencia de \textit{Evidence}, \textit{Originator}, \textit{Further Reading} y \textit{Basis} en este test preliminar respondió a la intención de maximizar la proporción de instancias de oro humano (etiquetadas por personas) en la evaluación de la línea base, reservando las instancias minadas para entrenamiento. El conjunto ampliado y el siguiente ciclo experimental resuelven este sesgo, incorporando las nueve clases en todas las particiones.

\section{Descripción del proceso de construcción de los modelos o solución}
\subsection{Diseño común y recuperación de evidencia}
El diseño compara tres familias bajo las entradas A/B, con la misma población y selección mediante validación (Figura~\ref{fig:flujo}). Se mantienen separados dos usos de SciBERT: un codificador para recuperar fragmentos, cuyos resultados quedan precomputados, y un clasificador ajustado para predecir funciones. Los experimentos aquí descritos evalúan clasificación sobre esos fragmentos; no vuelven a entrenar ni validan directamente el recuperador.

La segmentación respeta límites de sección y acumula oraciones contiguas. Frente al máximo de 300 palabras contemplado inicialmente, el diseño adoptó un umbral de unas 110 palabras para reducir la presión sobre la ventana de 512 tokens. Una oración individual que excede ese umbral se conserva íntegra, por lo que no es una cota absoluta. Reducir palabras tampoco garantiza una cantidad fija de tokens WordPiece.

Para recuperar, SciBERT codifica exclusivamente el contexto citante y el texto de cada chunk, sin prefijar encabezados retóricos. El propósito es evitar que la coincidencia entre nombres de sección dirija la búsqueda. Se promedian los estados ocultos de tokens válidos mediante la máscara de atención, se normaliza el vector en norma L2 y se calcula similitud coseno. La búsqueda se restringe a los fragmentos del documento citado objetivo, no a todo el corpus:
\[
 e(T)=\frac{\sum_i m_i h_i}{\sum_i m_i},\qquad
 s(q,c)=\frac{e(q)^\top e(c)}{\|e(q)\|_2\|e(c)\|_2}.
\]
Se ordenan los puntajes y se conservan hasta tres fragmentos, sus identificadores, secciones y posiciones. El almacenamiento del texto recuperado permite que distintas familias consuman la misma evidencia y evita recalcularla en cada entrenamiento. La similitud semántica es un criterio de selección, no una prueba de pertinencia argumentativa.

\begin{figure*}[t]
\centering
\begin{tikzpicture}[font=\small,box/.style={draw=blue!45!black,rounded corners=2pt,align=center,fill=blue!3,minimum height=0.85cm,text width=3.5cm},arr/.style={-{Stealth[length=2mm]},thick,draw=blue!50!black}]
\node[box,text width=11.6cm] (data) at (0,0) {Mismas instancias y particiones · contexto citante + referencia objetivo};
\node[box,text width=5.2cm] (a) at (-3.2,-1.4) {A: contexto + título + abstract};
\node[box,text width=5.2cm] (b) at (3.2,-1.4) {B: contexto + hasta 3 chunks\\recuperación dentro del citado};
\node[box] (tf) at (-4.2,-3.1) {TF-IDF + regresión logística\\baseline léxico};
\node[box] (sb) at (0,-3.1) {SciBERT ajustado\\control adicional: solo contexto};
\node[box] (qw) at (4.2,-3.1) {Qwen + LoRA\\ajuste de bajo rango};
\node[box,text width=11.6cm] (res) at (0,-4.65) {Función de cita · selección en dev · evaluación común en test\\Macro-F1 y cobertura por clase; comparación pareada A/B};
\draw[arr] (data.south) -- ++(0,-0.25) -| (a.north);
\draw[arr] (data.south) -- ++(0,-0.25) -| (b.north);
\foreach \x in {a,b} {\foreach \y in {tf,sb,qw} {\draw[arr] (\x.south) -- (\y.north);}}
\foreach \x in {tf,sb,qw} {\draw[arr] (\x.south) -- (\x.south |- res.north);}
\end{tikzpicture}
\caption{Diseño comparativo reproducible, dibujado en LaTeX/TikZ. Las nueve etiquetas de los experimentos iniciales difieren de las nueve funciones objetivo del conjunto ampliado; ambas versiones deben evaluarse por separado.}\label{fig:flujo}
\end{figure*}

\subsection{Baseline léxico: TF-IDF y regresión logística}
Se construyó una referencia léxica para medir cuánto puede resolverse a partir de palabras y expresiones del texto antes de introducir un modelo contextual. TF-IDF transforma cada entrada en un vector: pondera la frecuencia de los términos dentro del ejemplo y reduce el peso de los que aparecen en muchos documentos. Se emplean frecuencia sublineal, normalización L2 y un máximo de 100.000 términos. El vocabulario y los pesos IDF se aprenden exclusivamente en train.

La regresión logística aprende a asociar esos vectores con las funciones de cita. Se usan pesos de clase balanceados para dar mayor importancia relativa a las clases poco frecuentes, el optimizador \texttt{lbfgs}, un máximo de 2.000 iteraciones y semilla 42. Para A se concatena contexto, título y resumen; para B, contexto y fragmentos por rango. En ambas entradas se eliminan las marcas HTML de cita conservando su texto, sin añadir encabezados de sección ni limitar la longitud por tokens.

Las configuraciones r1, r2 y r3 se expresan mediante la tupla $(n_{\max},\mathrm{min\_df},C)$: $n_{\max}$ es la longitud máxima de los n-gramas (1: palabras; 2: palabras y pares consecutivos), \texttt{min\_df} es el mínimo de ejemplos de entrenamiento en que debe aparecer un término y $C$ es el inverso de la intensidad de regularización (un valor mayor penaliza menos los coeficientes). La Tabla~\ref{tab:tfconfig} hace explícitas las tres combinaciones.

\begin{table}[t]
\centering\small
\caption{Configuraciones TF-IDF y regresión logística, ejecutadas tanto para A como para B.}\label{tab:tfconfig}
\begin{tabular}{lcrr}\toprule
Configuración & N-gramas & \texttt{min\_df} & $C$\\\midrule
r1 & 1 & 1 & 0,5\\
r2 & 1--2 & 2 & 1,0\\
r3 & 1--2 & 3 & 2,0\\\bottomrule
\end{tabular}
\end{table}

Se entrenan seis modelos: tres configuraciones para cada entrada. La selección usa el mayor Macro-F1-9 en dev, es decir, el promedio del F1 de las nueve etiquetas, con el mismo peso para cada clase. La configuración r3 obtiene el mejor resultado de validación en ambas entradas; únicamente esos ganadores se evalúan en test, sin reentrenar con dev. Las combinaciones exploran conjuntamente vocabulario y regularización, por lo que se interpretan como configuraciones completas.

\subsection{Ajuste de SciBERT y control de contexto}
Se utiliza \texttt{allenai/scibert\_scivocab\_uncased}, un encoder preentrenado para texto científico \cite{Beltagy2019}, con una cabeza de clasificación de nueve salidas. A diferencia del baseline, la entrada incluye la sección del citante cuando está disponible; B añade también la sección de cada chunk. La variante \emph{solo contexto} elimina la información del citado, conservando el encuadre del citante, para medir si A/B aportan más que ese control.

La tokenización forma pares de secuencias: contexto y evidencia del citado. Se limita a 512 tokens con \texttt{longest\_first}, recortando la secuencia más larga; no se reserva una cuota uniforme para cada chunk. Este detalle es relevante porque B concatena los fragmentos por rango y los últimos pueden perder presencia. El modelo se ajusta con entropía cruzada ponderada por frecuencias de train, optimizador AdamW y planificación lineal de la tasa de aprendizaje.

\begin{table}[t]
\centering\small
\caption{Configuración de hiperparámetros para el ajuste de SciBERT.}\label{tab:hp}
\begin{tabularx}{\columnwidth}{@{}lY@{}}\toprule
Parámetro & Valor\\\midrule
Semillas & 42, 43, 44\\
Tasa de aprendizaje & $2\times10^{-5}$\\
Épocas máximas / paciencia & 5 / 2\\
Lote / acumulación & 8 / 2 (efectivo nominal: 16)\\
Lote de evaluación & 32\\
Warmup / weight decay & 0,10 / 0,01\\
Norma máxima de gradiente & 1,0\\
Longitud / selección & 512 / Macro-F1-9 en dev\\\bottomrule
\end{tabularx}
\end{table}

Para cada semilla y entrada se conserva el estado con mayor Macro-F1-9 en dev y se evalúa ese estado en test. Se ejecutan tres semillas para cada una de las tres entradas, sin escoger la mejor semilla por test. Las cifras agregadas son medias de métricas de corridas independientes, no métricas de un ensamble. El notebook también reejecuta el baseline TF-IDF sobre sus mismas particiones y reproduce los resultados iniciales, fortaleciendo la comparación interna.

El registro experimental conserva versiones, hashes de las particiones, hiperparámetros y predicciones por instancia. Esto permite reconstruir tablas y contrastes sin usar identificadores de corrida como contenido central del informe. Para reproducir la comparación deben preservarse además el orden de las etiquetas, la política de truncamiento y las diferencias de preprocesamiento: los contrastes entre familias no aíslan exclusivamente el efecto del encoder.

\subsection{Ajuste de Qwen+LoRA}
La siguiente incorporación es un modelo Qwen adaptado con LoRA, que introduce actualizaciones de bajo rango en el ajuste del modelo \cite{lora}. Su inclusión constituye una evolución del alcance: la propuesta inicial dejaba LoRA fuera del núcleo comparativo. La comparación actual lo incorpora como tercera familia, manteniendo la pregunta A/B.

\textbf{Formulación.} Se ajusta un LLM \emph{decoder} como clasificador de las 9 funciones de cita. El modelo no genera la etiqueta: una cabeza lineal ($d\to 9$) lee el estado oculto del último token y da un puntaje por clase, así que la salida siempre es una clase válida. En \cite{yousefiramandi}, con una sola etiqueta por ejemplo, esta formulación rindió mejor que la generativa basada en instrucciones. Se usan las versiones base de Qwen3-4B y Qwen3-8B, sin \emph{instruction tuning} \cite{qwen3}.

\textbf{Entrada.} El \emph{prompt} es una plantilla fija en inglés, sin instrucciones ni lista de clases, que termina siempre en la misma pregunta; su último token es el que lee la cabeza:
\begin{alltt}\footnotesize
Citing section: <sección retórica del citante>
Citation context: <contexto de la cita>
Cited paper: <título>. <abstract>   \(\leftarrow\) configuración A
Cited paper excerpts:               \(\leftarrow\) configuración B
[<encabezado de sección>] <fragmento 1>
[<encabezado de sección>] <fragmento 2>
[<encabezado de sección>] <fragmento 3>
Citation function:
\end{alltt}
A usa el título y el abstract del artículo citado; B, los tres fragmentos de su texto más similares al contexto, recuperados con SciBERT al construir el conjunto de datos. El texto no se normaliza. A diferencia de SciBERT, un LLM no está sujeto al truncamiento de 512 tokens; en este caso, y para compararlo mejor con SciBERT, el límite se fijó en 1.024 tokens. Solo se recorta el texto del citado; pasa en 17 instancias de entrenamiento de B (0,8\%) y en ninguna de validación ni de prueba.

\textbf{Adaptador.} LoRA \cite{lora} congela cada matriz de pesos $W\in\mathbb{R}^{d_\mathrm{out}\times d_\mathrm{in}}$ y aprende una corrección de rango bajo $\Delta W=(\alpha/r)\,BA$, con $A\in\mathbb{R}^{r\times d_\mathrm{in}}$ y $B\in\mathbb{R}^{d_\mathrm{out}\times r}$; como $B$ empieza en cero, el entrenamiento parte del modelo base. Se usó $r=16$, $\alpha=32$ (escala $\alpha/r=2$) y \emph{dropout} de 0,05 en la rama LoRA. El adaptador va en las siete proyecciones lineales de los 36 bloques ($W_q$, $W_k$, $W_v$ y $W_o$ en la atención; $W_\mathrm{gate}$, $W_\mathrm{up}$ y $W_\mathrm{down}$ en la MLP), 252 matrices: adaptar todas las capas lineales, y no solo la atención, es lo que acerca LoRA al ajuste completo \cite{schulman,qlora}. También se entrena la cabeza, que es nueva ($9d$ parámetros); la base queda congelada en bf16.

Cada matriz adaptada suma $r(d_\mathrm{in}+d_\mathrm{out})$ parámetros, así que se entrenan $36\,r\sum(d_\mathrm{in}+d_\mathrm{out})+9d$: 33,05\,M en el 4B ($d$ = 2.560, $d_\mathrm{ff}$ = 9.728), el 0,82\% del modelo, y 43,68\,M en el 8B ($d$ = 4.096, $d_\mathrm{ff}$ = 12.288), el 0,57\%. El porcentaje baja al crecer el modelo porque LoRA crece de forma lineal con las dimensiones, y la base, de forma cuadrática.

\textbf{Entrenamiento.} Entropía cruzada sin pesos por clase, AdamW sin \emph{weight decay} y recorte del gradiente en norma 1,0. El lote efectivo es de 16 ($4\times4$ de acumulación): 141 pasos por época sobre 2.255 ejemplos, con un máximo de 8 épocas (1.128 pasos). La tasa sigue un planificador WSD \cite{minicpm,hagele}: calentamiento lineal en el 5\% de los pasos, tasa constante y descenso hasta 0 con forma $1-\sqrt{\cdot}$ desde el paso 903 (el 20\% final). Se entrena en bf16 en una GPU NVIDIA RTX PRO 6000 de 96\,GB (Colab G4): cada corrida tomó de 11 a 19 min con el 4B y de 12 a 36 min con el 8B, con una memoria pico de 33--38\,GB y 46--53\,GB.

\textbf{Selección.} Se evalúa la validación dos veces por época; tras 5 evaluaciones seguidas sin mejora (2,5 épocas), el entrenamiento se detiene y se restauran los pesos del mejor punto. El criterio es el F1 macro de las 7 clases presentes en validación, que ordena igual que el Macro-F1-9. La prueba se evalúa una sola vez. Se reportan el Macro-F1-9 (las 9 clases; las ausentes cuentan 0) y el Macro-F1-5 (las 5 de prueba), con intervalo bootstrap del 95\%; en prueba, Macro-F1-9 = $5/9\cdot{}$Macro-F1-5.

\textbf{Corridas.} Se planeó una grilla de $2\times2\times3=12$ corridas: modelo (4B y 8B), entrada (A y B) y tasa de aprendizaje ($5\times10^{-5}$, $10^{-4}$ y $2\times10^{-4}$), que se barre en cada combinación porque la tasa óptima cambia con la configuración \cite{lee}. Lo demás queda fijo, con semilla 42. Se completaron 11; falta 8B-B con $2\times10^{-4}$ (Tabla~\ref{tab:qwen}).


\begin{table*}[t]
\centering\small
\caption{Resultados de test sobre 3.232 instancias particionadas. Macro-9 usa las nueve etiquetas iniciales; Macro-5 solo las cinco presentes en test. SciBERT: media de tres semillas. Qwen+LoRA: corrida con mayor Macro-F1-9 de validación en cada entrada, con una semilla.}\label{tab:results}
\begin{tabular}{llrrrr}\toprule
Modelo & Entrada & Macro-F1-9 & Macro-F1-5 & F1 ponderado & Exactitud\\\midrule
Referencia mayoritaria & Clase más frecuente de train & 0,0521 & 0,0939 & -- & 0,3066\\
TF-IDF + regresión logística & A & 0,1663 & 0,2993 & 0,4049 & 0,3745\\
TF-IDF + regresión logística & B & 0,1286 & 0,2315 & 0,3413 & 0,3045\\\midrule
SciBERT & Solo contexto & 0,2649 & 0,4767 & 0,5982 & 0,5988\\
SciBERT & A & 0,2648 & 0,4766 & 0,5935 & 0,5823\\
SciBERT & B & 0,2585 & 0,4652 & 0,5814 & 0,5768\\\midrule
Qwen3-4B-Base + LoRA & A & 0,2626 & 0,4726 & \textbf{0,5934} & 0,5823\\
Qwen3-8B-Base + LoRA & B & \textbf{0,2672} & \textbf{0,4809} & 0,5812 & \textbf{0,5967}\\\bottomrule
\end{tabular}
\end{table*}

\subsection{Métricas y contraste pareado}
Para cada clase, la \textbf{precisión} indica qué proporción de sus predicciones es correcta, y la \textbf{exhaustividad} (recall) qué proporción de sus ejemplos reales se recupera. El \textbf{F1} combina ambas mediante su media armónica:
\[
 F1_k=\frac{2P_kR_k}{P_k+R_k},\qquad
 \mathrm{Macro\text{-}F1}_{K}=\frac{1}{K}\sum_{k=1}^{K}F1_k.
\]
El \textbf{Macro-F1-9} promedia las nueve etiquetas del conjunto inicial y es el criterio común de selección en dev. El \textbf{Macro-F1-5} promedia únicamente las cinco clases con ejemplos en test, facilitando la lectura del desempeño sobre la cobertura disponible. Se fija en cero el F1 cuando su cálculo es indefinido. El \textbf{F1 ponderado} promedia los F1 según el número de ejemplos reales de cada clase; la \textbf{exactitud} es la proporción de predicciones correctas sobre todos los ejemplos. Por ello, Macro-F1 destaca el comportamiento por función, mientras el ponderado y la exactitud reflejan más las clases frecuentes.

El notebook de SciBERT calcula 10.000 remuestreos bootstrap pareados para diferencias de Macro-F1-9, tanto por instancia como por artículo citante. Se prioriza el segundo en la interpretación porque las citas de un mismo documento no son independientes. En cada remuestreo se conserva el emparejamiento de predicciones y se promedian las métricas de las semillas. Los intervalos expresan incertidumbre sobre este test; no representan la variación de todas las posibles particiones ni validan la taxonomía nueva.

\section{Resultados obtenidos}
\subsection{Selección y resultados del baseline}
La Tabla~\ref{tab:tfdev} resume la búsqueda en validación. r3 alcanza el mayor Macro-F1-9 para A (0,2696) y B (0,2209), por lo que establece la referencia de cada estrategia. En test obtiene 0,1663 para A y 0,1286 para B, frente a 0,0521 de la referencia que siempre predice la clase más frecuente de train. El baseline aporta así una referencia superior a la predicción constante en Macro-F1 y permite medir el beneficio posterior del ajuste contextual. La diferencia entre dev y test se analiza junto con su distinta composición documental y cobertura de clases.

\begin{table}[t]
\centering\small
\caption{Selección del baseline en dev. Se destaca el ganador por entrada; todas las filas usan Macro-F1-9.}\label{tab:tfdev}
\begin{tabular}{lrrrr}\toprule
 & \multicolumn{2}{c}{Entrada A} & \multicolumn{2}{c}{Entrada B}\\
Config. & Macro-F1 & Exactitud & Macro-F1 & Exactitud\\\midrule
r1 & 0,2190 & 0,3136 & 0,1900 & 0,3279\\
r2 & 0,2503 & 0,3381 & 0,2113 & 0,3544\\
r3 & \textbf{0,2696} & 0,3849 & \textbf{0,2209} & 0,3727\\\bottomrule
\end{tabular}
\end{table}


\subsection{Comparación conjunta sobre el conjunto inicial}
La Tabla~\ref{tab:results} reúne los resultados sobre los mismos 486 ejemplos de test. TF-IDF queda como referencia compacta, mientras SciBERT permite contrastar representaciones y el aporte del citado. En SciBERT se presentan medias de tres semillas; la desviación estándar de Macro-F1-9 es 0,0020 para contexto, 0,0145 para A y 0,0042 para B.



TF-IDF favorece A, con una diferencia B$-$A de $-0{,}0377$ en Macro-F1-9. En SciBERT la diferencia media es menor: $-0{,}0063$. El control de contexto alcanza 0,2649, ligeramente por encima de A y B; por tanto, en este corte no se observa una mejora clara al añadir información del citado. La Figura~\ref{fig:bars} muestra la comparación sin representar los resultados pendientes como ceros.

\begin{figure}[t]
\centering
\begin{tikzpicture}
\begin{axis}[width=\columnwidth,height=5.2cm,ybar,bar width=12pt,ymin=0,ymax=0.32,ylabel={Macro-F1-9},symbolic x coords={TF-IDF,SciBERT,Qwen + LoRA},xtick=data,legend style={at={(0.5,1.02)},anchor=south,legend columns=2,draw=none},ymajorgrids=true,grid style={gray!20},enlarge x limits=0.25]
\addplot[fill=blue!60,draw=blue!70] coordinates {(TF-IDF,0.1663) (SciBERT,0.2648) (Qwen + LoRA,0.2626)};
\addplot[fill=orange!75,draw=orange!90!black] coordinates {(TF-IDF,0.1286) (SciBERT,0.2585) (Qwen + LoRA,0.2672)};
\legend{A,B}
\end{axis}
\end{tikzpicture}
\caption{Comparación A/B en Macro-F1-9. SciBERT: media de tres semillas; el control solo contexto obtiene 0,2649. Qwen+LoRA: mejor configuración según validación, con una semilla por entrada (A: Qwen3-4B-Base; B: Qwen3-8B-Base).}\label{fig:bars}
\end{figure}

\begin{table}[t]
\centering\small
\caption{Diferencias de Macro-F1-9 e IC 95\% mediante bootstrap pareado por artículo citante.}\label{tab:ci}
\setlength{\tabcolsep}{3pt}
\begin{tabular}{lrr}\toprule
Contraste SciBERT & Diferencia & IC 95\%\\\midrule
B $-$ A & $-0,0063$ & [$-0,0142$; $0,0033$]\\
A $-$ contexto & $-0,0001$ & [$-0,0181$; $0,0165$]\\
B $-$ contexto & $-0,0064$ & [$-0,0222$; $0,0082$]\\
A $-$ TF-IDF A & $0,0985$ & [$0,0602$; $0,1397$]\\
B $-$ TF-IDF B & $0,1298$ & [$0,0956$; $0,1746$]\\\bottomrule
\end{tabular}
\end{table}

Los intervalos por documento para SciBERT frente a TF-IDF quedan por encima de cero en A y B. En cambio, los contrastes entre A, B y contexto incluyen cero (Tabla~\ref{tab:ci}). Esto respalda una ventaja de las configuraciones de SciBERT evaluadas frente al baseline sobre este test, pero no una superioridad concluyente entre sus fuentes de información.

\subsection{Barrido de Qwen+LoRA}
En prueba, el Macro-F1-5 va de 0,452 a 0,516 (Macro-F1-9, de 0,251 a 0,287) y los intervalos del 95\% de todas las corridas se solapan. Ningún factor tiene un efecto claro: entre pares que solo difieren en ese factor, B supera a A por 0,009 en promedio y 8B a 4B por 0,015, ambos con signo cambiante, y la tasa no sigue una tendencia. Por validación se elige 4B-A con $10^{-4}$, que en prueba da 0,473 de Macro-F1-5 y 0,263 de Macro-F1-9. La mejor en prueba, 8B-B con $10^{-4}$ (0,516 y 0,287), no puede elegirse por prueba sin sesgar el resultado. El mejor punto llegó entre las épocas 2 y 6, siempre antes del descenso de la tasa. Por clase, Background, Uses y Comparison quedan entre 0,56 y 0,73 de F1 en prueba; Extends (0,09--0,38) y Motivation (0,00--0,27), con 20 y 29 ejemplos, pesan el 40\% del Macro-F1-5 y concentran la mayor parte de la variación entre corridas. Con una sola semilla, estas diferencias no se distinguen de la variación entre semillas \cite{dodge}.

\begin{table}[t]
\centering\footnotesize
\caption{Qwen+LoRA: Macro-F1-9 de validación / Macro-F1-5 de prueba. En negrita, la corrida elegida por validación.}\label{tab:qwen}
\setlength{\tabcolsep}{3pt}
\begin{tabular}{lccc}\toprule
Modelo, entrada & $5\times10^{-5}$ & $10^{-4}$ & $2\times10^{-4}$\\\midrule
Qwen3-4B-Base, A & 0,413/0,462 & \textbf{0,492/0,473} & 0,401/0,488\\
Qwen3-4B-Base, B & 0,414/0,454 & 0,420/0,455 & 0,408/0,503\\
Qwen3-8B-Base, A & 0,428/0,491 & 0,408/0,452 & 0,368/0,467\\
Qwen3-8B-Base, B & 0,491/0,481 & 0,476/0,516 & --\\\bottomrule
\end{tabular}
\end{table}

\subsection{Cobertura, errores y truncamiento}
Las funciones frecuentes tienen mejor desempeño que las minoritarias. Para SciBERT A, el F1 medio es 0,6834 en Background, 0,6027 en Uses y 0,6159 en Comparison, frente a 0,2811 en Extends y 0,1999 en Motivation. En el baseline A esos dos últimos valores son 0,0789 y 0,1250. La mejora global no elimina la dificultad de las funciones con poco soporte, y nada permite inferir el desempeño de las cuatro clases ausentes.

La auditoría de longitud del notebook muestra que B excede 512 tokens en 47,4\% de train, 38,1\% de dev y 50,2\% de test. En A las tasas son 5,5\%, 0,2\% y 2,3\%. La medición permite identificar un ajuste concreto para el siguiente ciclo: distribuir el presupuesto de tokens entre los fragmentos. En la configuración actual, B utiliza la evidencia recuperada que permanece después del recorte a 512 tokens.

\section{Discusión de resultados y retos metodológicos}
\subsection{Qué permiten concluir los experimentos}
El avance establece un protocolo comparable y muestra que SciBERT mejora el baseline léxico, pero no confirma la hipótesis de que los fragmentos locales superen al título y resumen. Tampoco demuestra que la evidencia del citado sea innecesaria: el control de contexto compite con A/B en un test pequeño, concentrado en 16 documentos y limitado a cinco funciones de origen MultiCite. Es posible que muchas decisiones se resuelvan mediante señales del citante, pero esa explicación debe contrastarse en una evaluación con mayor cobertura.

La diferencia entre familias combina arquitectura, preentrenamiento, ajuste y preprocesamiento. SciBERT incorpora secciones y opera con un presupuesto de tokens; TF-IDF no. Por tanto, la comparación describe sistemas completos bajo una población común, no un experimento que varíe únicamente el encoder. Para aislar mecanismos conviene mantener constante el preprocesamiento y añadir ablaciones explícitas de secciones, longitud y cantidad de fragmentos.

La paridad A/B protege la comparación frente a cambios de ejemplos, pero restringe su dominio: solo incluye citas cuyo documento citado pudo recuperarse. Los trabajos sin cuerpo accesible pueden diferir en antigüedad, disciplina o forma de publicación. El aprovechamiento reducido del embudo inicial es así un resultado metodológico que debe reportarse, no solo un problema de almacenamiento. Recuperar más documentos amplía cobertura, pero no garantiza que desaparezca ese sesgo de disponibilidad.

\subsection{Calidad de etiquetas y modularidad del pipeline}
La versión ampliada refina la representación de las funciones: Basis se distingue de Extends y conserva su sentido de fundamento intelectual, mientras Improvement reúne las extensiones y modificaciones. Las reglas operativas hacen explícita esta decisión y permiten concentrar la revisión en los límites entre categorías. Motivation puede expresar necesidad de métodos o datos sin formular explícitamente un vacío; la absorción de Future Work en Gap también requiere examinar el contexto. La validación empírica debe focalizarse en estos límites y en los casos compuestos, sin deducir certeza semántica a partir de la marca formal del procedimiento de homologación.

Seleccionar la primera etiqueta permite una línea base reproducible de clasificación disyunta, pero no sustituye la adjudicación de la función discursiva dominante. Por ello, se propone comparar ese criterio con una revisión contextual y, como análisis de sensibilidad, con la exclusión de casos con mapeo múltiple (\textit{partial}). La preservación de metadatos exhaustivos de procedencia en el esquema del dataset resulta decisiva: al registrar la fuente de origen, la etiqueta original y el método de asignación, se habilita una auditoría sistemática con humanos en el bucle (\textit{human-in-the-loop}), permitiendo aislar con facilidad subconjuntos específicos de instancias para su posterior certificación.

La arquitectura del pipeline en capas desacopladas demostró una alta flexibilidad ante refinamientos metodológicos. Ante la evidencia experimental de que las entradas de la Configuración B excedían frecuentemente el límite de 512 tokens en SciBERT, la segmentación del texto citado se recalibró de un umbral inicial de 300 palabras a uno de aproximadamente 110 palabras. Este cambio se incorporó de forma inmediata mediante la parametrización de la capa de procesamiento de fragmentos, sin necesidad de re-descargar los datos crudos ni alterar los documentos estructurados previamente.

\section{Próximos pasos y plan de expansión}
\subsection{Siguiente ciclo: evidencia, cobertura y modelamiento}
\textbf{Control del presupuesto de evidencia y truncamiento.} Se evaluará la pertinencia de los fragmentos recuperados y se medirá qué proporción sobrevive a la tokenización de 512 tokens. En los LLM, el límite de tokens pasará a ser el máximo de contexto del modelo. Un presupuesto explícito por fragmento o una clasificación desagregada con combinación de puntajes permitirán comprobar si el recorte mecánico explica parte de la paridad observada entre A y B. Estas alternativas se seleccionarán en validación; los puntajes de similitud semántica orientan la selección, pero no reemplazan métricas de recuperación con referencias verificadas.

\textbf{Expansión y balanceo hacia 18.000 pares.} Se avanzará hacia la consolidación del corpus balanceado de 18.000 instancias pareadas (2.000 por función). La ampliación prioriza Gap, Improvement, Originator y Basis, cuyos déficits son mayores frente a clases saturadas como Further Reading o abundantes como Background. Se reportarán rigurosamente las tasas de admisión y descarte por fuente y función, distinguiendo la ausencia de texto completo, la ambigüedad en la resolución de identidad y la duplicación documental.

\textbf{Supervisión débil mediante Inferencia de Lenguaje Natural (NLI).} Para subsanar el déficit de las clases de bajo recall léxico sin incurrir en costos inviables de anotación manual, se integró supervisión débil basada en \texttt{DeBERTa-v3-base} (86M parámetros, preentrenado en NLI). A partir de premisas contextuales y de hipótesis discursivas derivadas de las definiciones canónicas \cite{Lyu2021}, el modelo predice la implicación textual de cada función. El protocolo se calibró empíricamente ante la sobre-abstención de los umbrales teóricos iniciales, fijando una probabilidad mínima $P(\text{top1}) \ge 0{,}70$, un margen de separación con la segunda hipótesis $\Delta \ge 0{,}15$ y temperatura $\tau = 0{,}7$. En un espacio de 9 hipótesis ($P_{\text{azar}} \approx 11{,}1\%$), $P \ge 0{,}70$ exige más de seis veces la probabilidad aleatoria y confiere al top-1 una ventaja superior a 18 veces el promedio de las hipótesis rivales ($\le 3{,}75\%$). Los resultados preliminares muestran que este mecanismo permite incorporar más de 5.000 instancias adicionales en clases complejas (\textit{Originator}, \textit{Basis}, \textit{Application} y \textit{Gap}) bajo control de saturación de cupos.

\textbf{Reevaluación comparativa de modelos.} Una vez congelada la versión ampliada del conjunto con particiones agrupadas anti-fuga y balance de clases, se reejecutarán TF-IDF, SciBERT y el ajuste de bajo rango con Qwen+LoRA \cite{lora} sobre las nueve funciones objetivo. No se compararán como equivalentes los resultados del conjunto inicial y del ampliado: varían la taxonomía, la procedencia, el volumen y la distribución de clases.

\textbf{Exploración de modelos.} Se explorarán versiones más modernas y/o posiblemente mejores que Qwen3, de la misma familia o de otras de pesos abiertos, ajustadas a la memoria de la GPU de 96\,GB; se evaluarán con el mismo protocolo y se elegirán en validación.

\textbf{Validación humana en la partición de prueba.} La propuesta original contemplaba preetiquetado asistido y validación humana de la prueba. Dado que el conjunto actual combina anotaciones heredadas, reglas léxicas y asignación por NLI, se mantiene la revisión humana del conjunto de prueba como trabajo pendiente. Gracias a la trazabilidad del pipeline, las instancias de prueba provenientes de asignación automática serán auditadas y validadas manualmente antes de la evaluación final, garantizando una partición de prueba de alta fidelidad en las nueve funciones canónicas.

{\small
\bibliographystyle{IEEEtran}
% \bibliography{referencias}
\bibliography{referencias_final}
}

\end{document}